% !TeX root = ../抽象代数笔记.tex
% 由主文档通过 \input 载入；请编译主文档。
\setcounter{chapter}{19} % 保留原书章号。
\chapter{模、子模与模同态}\label{ch:modules}

\begin{description}
  \item[学习目标] 用同一语言描述阿贝尔群与带线性变换的向量空间，理解域上的线性代数为什么不能直接照搬到一般环。
  \item[先修知识] 第\ref{ch:ideals}章及线性代数；与域论主线可独立衔接。
\end{description}

\section{从向量空间到模}
对环 $R$ 定义含幺左 $R$-模，明确标量作用的结合性、两种分配律和单位元条件。证明阿贝尔群与 $\mathbb Z$-模的对应，域上的模就是向量空间。对域 $F$ 上的向量空间 $V$ 及线性变换 $T:V\to V$，用 $f(x)\cdot v=f(T)v$ 定义 $F[x]$-模结构。

\section{子模、商模与同态定理}
定义子模（submodule）、模同态（module homomorphism）与商模（quotient module），建立核、像及模版本的同构定理。通过阿贝尔群与线性空间的例子比较三种同构定理的共同结构，并解释为何这里的商结构依靠加法子群对标量作用封闭。

\section{生成、自由模与基}
定义有限生成模、线性无关、基和自由模（free module）。用 $\mathbb Z/n\mathbb Z$（$n\geq2$）说明有限生成模未必自由。对一般环区分“由 $r$ 个元素生成”与“具有 $r$ 元基”；涉及基的基数唯一性时限定到本书后续使用的交换整环，不无条件推广向量空间结论。

\section{直和、挠元与正合列初步}
定义模的有限直和（direct sum），并解释无限直和要求元素只有有限个非零分量，与无限直积不同。对整环 $R$ 上的模定义挠元（torsion element）与挠子模，证明挠元全体构成子模。把短正合列（short exact sequence）作为选学语言引入，说明像等于核的含义及分裂时的直和解释。

\section{例题、练习与自检}
把 $\mathbb Z^2$ 的给定子群写成子模并计算商模，比较 $\mathbb Z$-模与 $\mathbb Q$-向量空间的生成问题。把具体矩阵转化为 $F[x]$-模，解释不变子空间就是子模。安排“无挠不一定自由”的一般反例，并明确主理想整环上的有限生成无挠模具有更强结论。

